\documentclass[../../main.tex]{subfiles}

\begin{document}

\section{参考文献}\label{section:ch11-参考文献}

T.-P. Liu 撰写了本章全部内容的初稿，其书 \cite{Lu} 是本学科的主要参考文献。

\textbf{\S11.1：} 浅水波方程的讨论见 Holden--Risebro \cite{H-R} 和 LeVeque \cite{LV}。P. Colella 协助了 \S11.1.2 中的计算。\cref{theorem:ch11-特征值光滑依赖} 的证明遵循 W. Han 的建议。更多内容见 Bressan \cite{Bs}、Courant--Friedrichs \cite{C-F}、Dafermos \cite{D}、Majda \cite{Mj}、Xiao--Zhang \cite{X-Z} 和 Zheng \cite{Zh}。

\textbf{\S11.2：} 见 Lax (Comm. Pure and Applied Math 10 (1957), 537--566)。

\textbf{\S11.3：} 例子遵循 Logan \cite{Lo}。关于使用 Riemann 不变量的更多内容见 Neu \cite{Ne1}。

\textbf{\S11.4：} 关于粘性行波的更多内容见 Smoller \cite{S}。\S11.4.3 中的唯一性定理归功于 S. Kru\v{z}kov (Mat. Sbornik 123 (1970), 228--255)。

\textbf{\S11.5：} 问题 3 取自 Kato \cite{K}。问题 4--6 基于 Holden--Risebro \cite{H-R}。问题 10 和 11 来自 Smoller \cite{S}，问题 13 来自 Dafermos \cite{D}。

\end{document}